\documentclass{article}

\usepackage{geometry} % This package allows the editing of the page layout
\usepackage{amsmath}  % This package allows the use of a large range
% of mathematical formula, commands, and symbols
\usepackage{graphicx}  % This package allows the importing of images
\usepackage{
  amsmath,      % Math Environments
  amssymb,      % Extended Symbols
  enumerate,        % Enumerate Environments
  graphicx,      % Include Images
  lastpage,      % Reference Lastpage
  multicol,      % Use Multi-columns
  multirow      % Use Multi-rows
}

\newcommand{\question}[2][]{
  \begin{flushleft}
    \textbf{Question #1}: \textit{#2}

\end{flushleft}}
\newcommand{\sol}{\textbf{Solution}:} %Use if you want a boldface solution line
\newcommand{\maketitletwo}[2][]{
  \begin{center}
    \Large{\textbf{Assignment #1}

    MATH4362} % Name of course here
    \vspace{5pt}

    \normalsize{N. Ohayon Rozanes% Your name here

    \today}        % Change to due date if preferred
    \vspace{15pt}

\end{center}}
\begin{document}
\maketitletwo[4]  % Optional argument is assignment number
%Keep a blank space between maketitletwo and \question[1]

\question[1]{}
$$
\begin{cases}
  u_{tt} = 3u_{xx} \\
  u(t,0) = u(t,\pi) = 0  \\
  u(0,x) = \sin^3(x) \\
  u_t(0,x) = 0
\end{cases}
$$

From the seperable solutions to the wave equation, we have that $w(t)
= \cos \omega t$ or $\sin\omega t$ and $v(x) = \cos \omega x/c$ or
$\sin \omega x/c$, conditioning on the boundary conditions, $u(t,0) =
u(t, \pi) = 0$ yields that it must be the case that $v(x) = \sin \omega x/c$

Thus we have that:

$$
u = \sum b_n \cos \left(\frac{\pi nct}{l}\right)\sin\left(\frac{\pi
nx}{l}\right) + d_n
\sin\left(\frac{\pi nct}{l}\right)\sin\left(\frac{\pi nx}{l}\right)
$$

Since $l = \pi$, we can simplify the solution a little to:

$$
u = \sum b_n \cos \left(nct\right)\sin\left(
nx\right) + d_n
\sin\left(nct\right)\sin\left(nx\right)
$$

Taking the partial derivative with respect to $t$:

\[
  u_t = \sum -b_nnc\sin(nct)\sin(nx) + d_nnc\cos(nct)\sin(nx)
\]

conditioning on $u_t(0, x) = 0$ we get that:
$$
u_t(0,x) = \sum d_nnc\sin(nx) = 0
$$

It follows that for all $n$, $d_n = 0$. So the solution becomes

$$
u = \sum b_n \cos \left(nct\right)\sin\left(nx\right)
$$

Conditioning on the initial displacement we have that:

$$
\sin^3(x) = \sum b_n \sin(nx)
$$

Since $\sin$ is odd, $\sin^2$ is even, and $\sin^3$ is again odd,
therefore we can use the Fourier Sine coefficients to find $b_n$:
\begin{align*}
  b_n &=  \frac{2}{\pi}\int_0^\pi \sin^3(x)\sin(nx) dx \\
  &=  \frac{2}{\pi}\int_0^\pi \sin^2(x)\sin(x)\sin(nx) dx \\
  &=  \frac{1}{\pi}\int_0^\pi (1-\cos(2x))\sin(x)\sin(nx) dx \\
  &=  \frac{1}{\pi}\int_0^\pi (1-\cos(2x))
  \left(\frac{\cos((n-1)x) - \cos((n+1)x)}{2}\right) dx \\
  &= \frac{1}{2\pi}\int_0^\pi \Big(\cos((n-1)x) - \cos((n+1)x) \\
  &\qquad - \cos(2x)\cos((n-1)x) + \cos(2x)\cos((n+1)x)\Big) dx
\end{align*}

Using $\cos A\cos B = \frac{1}{2}(\cos(A-B)+\cos(A+B))$ and
$\int_0^\pi \cos(kx)dx = 0$ for $k \ne 0$, the only nonzero cases are
$n=1$ and $n=3$. This gives
\[
  b_n =
  \begin{cases}
    \frac{3}{4}, & n=1 \\
    -\frac{1}{4}, & n=3 \\
    0, & \text{otherwise}.
  \end{cases}
\]

Thus the solution is
\[
  u(t,x) = \frac{3}{4}\cos(ct)\sin(x) - \frac{1}{4}\cos(3ct)\sin(3x),
\]
with $c=\sqrt{3}$.
\question[2]{}
$$
\begin{cases}
  u_{tt} = u_{xx} \\
  u_x(t, 0) = u_x(t, \pi ) = 0
\end{cases}
$$

Again, starting from the possible solutions to the seperable
equations, we have that:
\begin{align*}
  u(t,x) &= v(x)w(t)\\
  u_x(t,x) &= v'(x)w(t) \\
\end{align*}
So,
\[
  v'(0)w(t) = v'(\pi)w(t) \implies v'(0) = v'(\pi) = 0
\]

For the sine candidate $v(x)=\sin(\omega x/c)$, we have
\[
  v'(x)=\frac{\omega}{c}\cos\left(\frac{\omega x}{c}\right),
\]
so $v'(0)=\frac{\omega}{c}\neq 0$ in general, which violates the
Neumann condition $v'(0)=0$. Thus the sine branch does not work.

However, for the $\cos$ candidate solution for the $v(x) =
\cos(\omega x/c)$, we have that:
\[
  v'(x) = \frac{\omega}{c}\sin\left(\frac{\omega x}{c}\right)
\]
plugging into $v'(0) = v'(\pi)$
\[
  \frac{\omega}{c}\sin\left(0\right) =
  \frac{\omega}{c}\sin\left(\frac{\omega \pi}{c}\right)
\]
Which simplifies to:
\[
  \sin(0) = \sin(\frac{\omega\pi}{c}) = 0
\]

Which is true for $\frac{\omega}{c} \in \mathbb{N}$, thus $\omega = nc$.

For the case $c^2=1$ (so $c=1$), the separated solutions are
$v_n(x)=\cos(nx)$ with $n=0,1,2,\dots$. Notice that the Neumann
conditions do not restrict the solutions of $w(t)$ at all, so the
total solution is:

$$
u(t,x) = \sum b_n\cos(nt)\cos(nx) + d_n\sin(nt)\cos(nx)
$$

\question[3]{}
\begin{enumerate}[(a)]
  \item
    \[
      \begin{cases}
        u_{tt} = u_{xx} \\
        u(t,\pi) = u(t, -\pi) \\
        u_x(t, \pi) = u_x(t, -\pi) \\
        u(0, x) = f(x) \\
        u_t(0, x) = g(x)
      \end{cases}
    \]
  \item From the condition $u(t, -\pi) = u(t, \pi)$, then we have that:
    \[
      v(-\pi)w(t) = v(\pi)w(t) \implies v(-\pi) = v(\pi)
    \]
    For $v(x) = \cos(\omega x/c)$, since cosine is an even function,
    the above condition is immediatly satisfied. For the sine
    candidate solution to $v$, we have that
    $$
    \sin(\frac{\omega\pi}{c}) = \sin(-\frac{\omega\pi}{c}) =
    -\sin(\frac{\omega\pi}{c})
    $$
    Which can only happen at $\sin(\frac{\omega\pi}{c}) = 0$,
    equivalently, when $\omega\pi/c \in \mathbb{N}$ which yields that
    $\omega_n = nc$. The same restriction follows from the derivative
    condition $u_x(t,-\pi)=u_x(t,\pi)$: with $u_x=v'(x)w(t)$ we need
    $v'(-\pi)=v'(\pi)$. For $v(x)=\sin(\omega x/c)$ we get
    $v'(x)=\frac{\omega}{c}\cos(\omega x/c)$, so
    \[
      \cos\left(\frac{\omega\pi}{c}\right)
      = \cos\left(-\frac{\omega\pi}{c}\right),
    \]
    which is automatic, but for $v(x)=\cos(\omega x/c)$ we have
    $v'(x)=-\frac{\omega}{c}\sin(\omega x/c)$ and the condition
    \[
      \sin\left(\frac{\omega\pi}{c}\right)
      = -\sin\left(\frac{\omega\pi}{c}\right)
    \]
    again forces $\sin(\omega\pi/c)=0$, so $\omega=nc$. With $\omega=nc$, the
    separated factors satisfy
    \[
      v_n'' + n^2 v_n = 0, \qquad w_n'' + n^2 c^2 w_n = 0,
    \]
    so
    \[
      v_n(x) = c_n\sin(nx) + d_n\cos(nx), \qquad
      w_n(t) = a_n\cos(ntc) + b_n\sin(ntc).
    \]
    Thus we get the solution
    \begin{align*}
      u(t, x)
      &= \sum (a_n\cos(ntc) + b_n\sin(ntc))(c_n\sin(nx) + d_n\cos(nx)) \\
      &= \sum \Big(
        \big(a_n d_n \cos(ntc) + b_n d_n \sin(ntc)\big)\cos(nx) +
      \big(a_n c_n \cos(ntc) + b_n c_n \sin(ntc)\big)\sin(nx)\Big)
    \end{align*}
    and therefore, with $A_n=a_n d_n$, $B_n=a_n c_n$, $C_n=nc\,b_n d_n$,
    and $D_n=nc\,b_n c_n$,
    \begin{align*}
      f(x) &= u(0,x) = \sum \big(A_n \cos(nx) + B_n \sin(nx)\big), \\
      g(x) &= u_t(0,x) = \sum \big(C_n \cos(nx) + D_n \sin(nx)\big).
    \end{align*}
    Using the Fourier series of $f(x)$ and $g(x)$ on $[-\pi,\pi]$,
    \begin{align*}
      f(x) &= \frac{A_0}{2} + \sum_{n=1}^\infty
      \left(A_n\cos(nx)+B_n\sin(nx)\right), \\
      g(x) &= \frac{C_0}{2} + \sum_{n=1}^\infty
      \left(C_n\cos(nx)+D_n\sin(nx)\right),
    \end{align*}
    where
    \begin{align*}
      A_n &= \frac{1}{\pi}\int_{-\pi}^{\pi} f(x)\cos(nx)\,dx,
      & B_n &= \frac{1}{\pi}\int_{-\pi}^{\pi} f(x)\sin(nx)\,dx, \\
      C_n &= \frac{1}{\pi}\int_{-\pi}^{\pi} g(x)\cos(nx)\,dx,
      & D_n &= \frac{1}{\pi}\int_{-\pi}^{\pi} g(x)\sin(nx)\,dx.
    \end{align*}
  \item Holding $x$ constant we have that
    \[
      u(t)
      = \sum (a_n\cos(ntc) + b_n\sin(ntc))K
    \]
    inspecting a single term of the sum, we can solve for the period:
    \begin{align*}
      u_n(t) &= a_n\cos(ntc) + b_n\sin(ntc), \\
      u_n(t+p) &= a_n\cos(nc(t+p)) + b_n\sin(nc(t+p)).
    \end{align*}
    So we need $nc(t+p)=nct+2\pi k$, hence $p=2\pi k/(nc)$ and the
    period is $2\pi/c$.
  \item The solution will behave in a basically similar way,
    beginning as a peak and splitting into two waves travelling in
    opposite directions, however, instead of flipping, the two waves
    "meet" again and eventually form the shape of the initial
    displacement again, repeating.
\end{enumerate}
\end{document}
